\documentclass[aspectratio=169]{beamer}
\usetheme{metropolis}
\usepackage{graphicx}
\usepackage{tikz}
\usepackage{amsmath}
\usepackage{listings}
\usepackage{xcolor}
\usepackage{colortbl}
\usepackage{booktabs}
\usepackage{emoji}
\usetikzlibrary{shapes,arrows,positioning,calc}

% Code listing setup
\lstset{
  basicstyle=\ttfamily\small,
  keywordstyle=\color{blue},
  commentstyle=\color{gray},
  stringstyle=\color{red},
  showstringspaces=false,
  breaklines=true,
  frame=single,
  language=Python
}

\title{Context-Aware Models}
\subtitle{Weeks 5-6: From Sequence-to-Sequence to Transformers}
\author{LLM Course}
\date{\today}

\begin{document}

% ============================================================
% TITLE SLIDE
% ============================================================
\begin{frame}
\titlepage
\end{frame}

% ============================================================
% OVERVIEW
% ============================================================
\begin{frame}{Overview 📋}
\large
\begin{enumerate}
    \item 🔄 \textbf{Sequence-to-Sequence Models}: The Foundation
    \item ⚡ \textbf{Attention Mechanisms}: The Key Innovation
    \item 🤖 \textbf{The Transformer}: Architecture Deep Dive
    \item 🎯 \textbf{BERT \& Variants}: Masked Language Modeling
    \item 🧠 \textbf{Cognitive Neuroscience}: Prediction in Brains vs. Models
    \item ✍️ \textbf{Assignment 4}: Building Context-Aware Systems
\end{enumerate}

\vspace{1em}
\textit{Goal: Understand how modern models leverage context to understand language}
\end{frame}

% ============================================================
% THE CONTEXT PROBLEM
% ============================================================
\begin{frame}{The Context Problem 🤔}
\begin{center}
\LARGE
\textbf{Why do we need context-aware models?}
\end{center}

\pause
\vspace{1em}

\begin{exampleblock}{Example: The word "bank"}
\begin{enumerate}
    \item "I deposited money at the \textbf{bank}" (financial institution)
    \item "We sat by the river \textbf{bank}" (riverside)
    \item "The plane started to \textbf{bank} left" (tilt/turn)
\end{enumerate}
\end{exampleblock}

\vspace{1em}

\begin{columns}[T]
\begin{column}{0.48\textwidth}
\textbf{Static Embeddings (Word2Vec):}
\begin{itemize}
    \item One vector per word
    \item \alert{Can't distinguish meanings!}
    \item Context-independent
\end{itemize}
\end{column}

\pause

\begin{column}{0.48\textwidth}
\textbf{Context-Aware Models:}
\begin{itemize}
    \item Different representation per occurrence
    \item Uses surrounding context
    \item Handles polysemy naturally
\end{itemize}
\end{column}
\end{columns}
\end{frame}

% ============================================================
% SEQUENCE TO SEQUENCE INTRO
% ============================================================
\begin{frame}{Sequence-to-Sequence Models 🔄}
\textbf{The breakthrough for variable-length input/output problems}

\vspace{1em}

\textbf{Applications:}
\begin{itemize}
    \item Machine Translation: English → French
    \item Summarization: Long text → Short summary
    \item Question Answering: Question + Context → Answer
    \item Dialogue Systems: User input → System response
\end{itemize}

\vspace{1em}

\begin{center}
\begin{tikzpicture}[scale=0.8]
    % Input sequence
    \foreach \i/\word in {1/The, 2/cat, 3/sat} {
        \node[rectangle,draw,fill=blue!20] at (\i*1.5, 2) {\word};
    }
    \node at (0, 2) {Input:};

    % Arrow
    \draw[->, very thick] (5, 2) -- (6.5, 2);
    \node[rectangle,draw,fill=green!20] at (7.5, 2) {Model};
    \draw[->, very thick] (8.5, 2) -- (10, 2);

    % Output sequence
    \node at (10.5, 2) {Output:};
    \foreach \i/\word in {1/Le, 2/chat, 3/s'assit} {
        \node[rectangle,draw,fill=red!20] at (10.5+\i*1.5, 2) {\word};
    }
\end{tikzpicture}
\end{center}

\vspace{1em}
\small
\textit{Reference: Sutskever et al. (2014) - "Sequence to Sequence Learning with Neural Networks"}
\end{frame}

% ============================================================
% ENCODER DECODER ARCHITECTURE
% ============================================================
\begin{frame}[fragile]{Encoder-Decoder Architecture 🏗️}
\textbf{Two-part architecture: Encode then Decode}

\vspace{1em}

\begin{center}
\begin{tikzpicture}[
    node distance=0.8cm,
    rnn/.style={rectangle,draw,fill=blue!20,minimum width=1cm,minimum height=1cm},
    decoder/.style={rectangle,draw,fill=red!20,minimum width=1cm,minimum height=1cm}
]
    % Encoder
    \node[rnn] (e1) {RNN};
    \node[rnn,right=of e1] (e2) {RNN};
    \node[rnn,right=of e2] (e3) {RNN};

    % Input words
    \node[below=of e1] {The};
    \node[below=of e2] {cat};
    \node[below=of e3] {sat};

    % Context vector
    \node[rectangle,draw,fill=green!20,minimum width=1.5cm,minimum height=1cm,right=of e3,align=center] (ctx) {Context\\Vector};

    % Decoder
    \node[decoder,right=of ctx] (d1) {RNN};
    \node[decoder,right=of d1] (d2) {RNN};
    \node[decoder,right=of d2] (d3) {RNN};

    % Output words
    \node[above=of d1] {Le};
    \node[above=of d2] {chat};
    \node[above=of d3] {s'assit};

    % Arrows
    \draw[->] (e1) -- (e2);
    \draw[->] (e2) -- (e3);
    \draw[->] (e3) -- (ctx);
    \draw[->] (ctx) -- (d1);
    \draw[->] (d1) -- (d2);
    \draw[->] (d2) -- (d3);

    % Labels
    \node[above=2cm of e2,blue] {\textbf{Encoder}};
    \node[above=2cm of d2,red] {\textbf{Decoder}};
\end{tikzpicture}
\end{center}

\vspace{1em}

\begin{columns}[T]
\begin{column}{0.48\textwidth}
\textbf{Encoder:}
\begin{itemize}
    \item Reads input sequence
    \item Compresses to fixed-size vector
    \item Captures semantic meaning
\end{itemize}
\end{column}

\begin{column}{0.48\textwidth}
\textbf{Decoder:}
\begin{itemize}
    \item Starts from context vector
    \item Generates output sequence
    \item One token at a time
\end{itemize}
\end{column}
\end{columns}
\end{frame}

% ============================================================
% SEQ2SEQ PROBLEM
% ============================================================
\begin{frame}{The Seq2Seq Bottleneck Problem ⚠️}
\textbf{Challenge: All information compressed into single vector!}

\vspace{1em}

\begin{alertblock}{The Problem}
\begin{itemize}
    \item Long sequences → information loss
    \item Fixed-size context vector is a bottleneck
    \item Early tokens forgotten by the time we reach the end
    \item Performance degrades with sequence length
\end{itemize}
\end{alertblock}

\vspace{1em}

\begin{center}
\begin{tikzpicture}[scale=0.9]
    % Long input
    \foreach \i in {1,...,8} {
        \node[rectangle,draw,fill=blue!20,minimum size=0.5cm] at (\i*0.7, 2) {};
    }
    \node at (0, 2) {Input:};

    % Bottleneck
    \draw[->] (6.5, 2) -- (7.5, 2);
    \node[rectangle,draw,fill=orange!50,minimum width=0.8cm,minimum height=1.5cm] at (8.5, 2) {};
    \node[below=0.1cm of {(8.5,2)}] {\alert{Bottleneck!}};
    \draw[->] (9.5, 2) -- (10.5, 2);

    % Output
    \foreach \i in {1,...,8} {
        \node[rectangle,draw,fill=red!20,minimum size=0.5cm] at (10.5+\i*0.7, 2) {};
    }
\end{tikzpicture}
\end{center}

\vspace{1em}

\begin{center}
\LARGE
\textbf{Solution: Attention! ⚡}
\end{center}
\end{frame}

% ============================================================
% ATTENTION MECHANISM
% ============================================================
\begin{frame}{Attention Mechanism: The Big Idea ⚡}
\textbf{Instead of compressing everything into one vector...}

\vspace{0.5em}

\textbf{Let the decoder look at all encoder hidden states!}

\vspace{1em}

\begin{center}
\begin{tikzpicture}[
    node distance=0.5cm,
    enc/.style={rectangle,draw,fill=blue!20,minimum width=0.8cm,minimum height=0.8cm},
    dec/.style={rectangle,draw,fill=red!20,minimum width=0.8cm,minimum height=0.8cm}
]
    % Encoder states
    \node[enc] (e1) at (0, 0) {$h_1$};
    \node[enc] (e2) at (1.2, 0) {$h_2$};
    \node[enc] (e3) at (2.4, 0) {$h_3$};
    \node[enc] (e4) at (3.6, 0) {$h_4$};

    \node[left=0.3cm of e1] {Encoder:};

    % Decoder state
    \node[dec] (d1) at (1.8, 3) {$s_t$};
    \node[left=0.3cm of d1] {Decoder:};

    % Attention arrows
    \draw[->, thick, blue] (e1) -- (d1) node[midway,left] {\tiny 0.1};
    \draw[->, thick, orange, line width=2pt] (e2) -- (d1) node[midway,left] {\tiny \textbf{0.6}};
    \draw[->, thick, orange, line width=2pt] (e3) -- (d1) node[midway,right] {\tiny \textbf{0.25}};
    \draw[->, thick, blue] (e4) -- (d1) node[midway,right] {\tiny 0.05};

    \node[below=1.5cm of e2] {Input: "The cat sat down"};
    \node[above=0.5cm of d1] {Generating: "Le \textbf{chat}"};
    \node[right=2cm of d1] {\alert{Attention weights sum to 1.0}};
\end{tikzpicture}
\end{center}

\vspace{1em}

\textbf{Key Insight:} When generating "chat" (cat), pay more attention to "cat" in the input!

\vspace{0.5em}
\small
\textit{Reference: Bahdanau et al. (2015) - "Neural Machine Translation by Jointly Learning to Align and Translate"}
\end{frame}

% ============================================================
% ATTENTION COMPUTATION
% ============================================================
\begin{frame}{How Attention Works: Step by Step 🔍}
\textbf{Computing attention weights:}

\vspace{1em}

\begin{enumerate}
    \item \textbf{Score}: How relevant is each encoder state to current decoder state?
    \begin{align*}
    e_{t,i} = \text{score}(s_t, h_i) = s_t^T W_a h_i
    \end{align*}

    \item \textbf{Normalize}: Convert scores to probabilities (softmax)
    \begin{align*}
    \alpha_{t,i} = \frac{\exp(e_{t,i})}{\sum_{j=1}^{T} \exp(e_{t,j})}
    \end{align*}

    \item \textbf{Context}: Weighted sum of encoder states
    \begin{align*}
    c_t = \sum_{i=1}^{T} \alpha_{t,i} h_i
    \end{align*}

    \item \textbf{Decode}: Use context vector along with decoder state
    \begin{align*}
    s_{t+1} = f(s_t, c_t, y_t)
    \end{align*}
\end{enumerate}

\vspace{0.5em}

\textbf{Result:} Decoder dynamically focuses on different parts of input!
\end{frame}

% ============================================================
% ATTENTION VISUALIZATION
% ============================================================
\begin{frame}{Attention Visualization 👁️}
\textbf{Example: English → French translation with attention weights}

\vspace{1em}

\begin{center}
\begin{tabular}{c|cccc}
\textbf{Output} & \multicolumn{4}{c}{\textbf{Input Words}} \\
\textbf{(French)} & The & agreement & on & European \\
\hline
L' & \cellcolor{orange!80}0.7 & 0.1 & 0.1 & 0.1 \\
accord & 0.1 & \cellcolor{orange!80}0.8 & 0.05 & 0.05 \\
sur & 0.05 & 0.05 & \cellcolor{orange!80}0.8 & 0.1 \\
l'européen & 0.05 & 0.05 & 0.1 & \cellcolor{orange!80}0.8 \\
\end{tabular}
\end{center}

\vspace{1em}

\textbf{Observations:}
\begin{itemize}
    \item Diagonal pattern for similar word order
    \item Model learns alignment automatically!
    \item No need for explicit word alignment annotations
    \item Can handle reordering (syntax differences between languages)
\end{itemize}

\vspace{0.5em}
\textit{Attention provides interpretability: we can see what the model is "looking at"}
\end{frame}

% ============================================================
% TRANSFORMER INTRO
% ============================================================
\begin{frame}{The Transformer Revolution 🤖}
\begin{center}
\LARGE
\textbf{"Attention Is All You Need"}
\end{center}

\vspace{1em}

\begin{columns}[T]
\begin{column}{0.48\textwidth}
\textbf{Before Transformers (2017):}
\begin{itemize}
    \item RNNs/LSTMs with attention
    \item Sequential processing
    \item Hard to parallelize
    \item Limited context window
    \item Slow training
\end{itemize}
\end{column}

\pause

\begin{column}{0.48\textwidth}
\textbf{After Transformers:}
\begin{itemize}
    \item \alert{Pure attention, no recurrence!}
    \item Parallel processing
    \item Scales to GPUs/TPUs
    \item Long-range dependencies
    \item Fast \& effective
\end{itemize}
\end{column}
\end{columns}

\vspace{1em}

\begin{block}{Key Innovation}
Replace sequential RNNs with \textbf{self-attention} layers that process all positions in parallel!
\end{block}

\vspace{0.5em}
\small
\textit{Reference: Vaswani et al. (2017) - "Attention Is All You Need"}
\end{frame}

% ============================================================
% TRANSFORMER ARCHITECTURE
% ============================================================
\begin{frame}[fragile]{Transformer Architecture Overview 🏗️}
\begin{center}
\begin{tikzpicture}[scale=0.65]
    % Encoder
    \draw[thick,blue] (0, 0) rectangle (3, 8);
    \node[blue] at (1.5, 8.5) {\textbf{Encoder}};

    % Encoder components
    \node[rectangle,draw,fill=blue!20,minimum width=2.5cm] at (1.5, 7) {Feed Forward};
    \node[rectangle,draw,fill=blue!30,minimum width=2.5cm,align=center] at (1.5, 5.5) {Multi-Head\\Attention};
    \node[rectangle,draw,fill=blue!20,minimum width=2.5cm] at (1.5, 4) {Feed Forward};
    \node[rectangle,draw,fill=blue!30,minimum width=2.5cm,align=center] at (1.5, 2.5) {Multi-Head\\Attention};
    \node[align=center] at (1.5, 1) {+ Positional\\Encoding};
    \node[below] at (1.5, 0) {Input};

    % Decoder
    \draw[thick,red] (5, 0) rectangle (8, 8);
    \node[red] at (6.5, 8.5) {\textbf{Decoder}};

    % Decoder components
    \node[rectangle,draw,fill=red!20,minimum width=2.5cm] at (6.5, 7) {Feed Forward};
    \node[rectangle,draw,fill=red!30,minimum width=2.5cm,align=center] at (6.5, 5.8) {Cross\\Attention};
    \node[rectangle,draw,fill=red!40,minimum width=2.5cm,align=center] at (6.5, 4.3) {Masked\\Self-Attention};
    \node[rectangle,draw,fill=red!20,minimum width=2.5cm] at (6.5, 2.8) {Feed Forward};
    \node[rectangle,draw,fill=red!30,minimum width=2.5cm,align=center] at (6.5, 1.3) {Masked\\Self-Attention};
    \node[below] at (6.5, 0) {Output};

    % Arrows
    \draw[->, thick] (3, 6) -- (5, 6) node[midway,above] {\tiny encoder states};
    \draw[->, thick] (3, 3.5) -- (5, 3.5);

    % Stack indicators
    \node[right=0.2cm of {(3,5.5)}] {\tiny Nx};
    \node[right=0.2cm of {(8,5.5)}] {\tiny Nx};
\end{tikzpicture}
\end{center}

\textbf{Key Components:}
\begin{itemize}
    \item \textbf{Multi-Head Self-Attention}: Relate all positions to each other
    \item \textbf{Feed-Forward Networks}: Transform representations
    \item \textbf{Residual Connections \& Layer Norm}: Training stability (not shown)
    \item \textbf{Positional Encoding}: Inject position information
\end{itemize}
\end{frame}

% ============================================================
% SELF ATTENTION
% ============================================================
\begin{frame}{Self-Attention: The Core Mechanism 🎯}
\textbf{Key idea: Each word attends to all other words in the sequence}

\vspace{1em}

\textbf{Three learned projections:}
\begin{itemize}
    \item \textbf{Query (Q)}: What am I looking for?
    \item \textbf{Key (K)}: What do I contain?
    \item \textbf{Value (V)}: What information do I have?
\end{itemize}

\vspace{1em}

\textbf{Computation:}
\begin{align*}
Q = XW_Q, \quad K = XW_K, \quad V = XW_V
\end{align*}
\begin{align*}
\text{Attention}(Q, K, V) = \text{softmax}\left(\frac{QK^T}{\sqrt{d_k}}\right)V
\end{align*}

\vspace{1em}

\begin{center}
\begin{tikzpicture}[scale=0.8]
    \node[rectangle,draw,fill=blue!20,minimum width=2cm,minimum height=0.8cm] (x) at (0, 0) {Input $X$};

    \node[rectangle,draw,fill=green!20,minimum width=1.5cm,minimum height=0.8cm] (q) at (-2, 2) {$Q$};
    \node[rectangle,draw,fill=orange!20,minimum width=1.5cm,minimum height=0.8cm] (k) at (0, 2) {$K$};
    \node[rectangle,draw,fill=purple!20,minimum width=1.5cm,minimum height=0.8cm] (v) at (2, 2) {$V$};

    \draw[->] (x) -- (q) node[midway,left] {$W_Q$};
    \draw[->] (x) -- (k) node[midway,left] {$W_K$};
    \draw[->] (x) -- (v) node[midway,right] {$W_V$};

    \node[rectangle,draw,fill=red!20,minimum width=2cm,minimum height=0.8cm] (attn) at (0, 4) {Attention};
    \draw[->] (q) -- (attn);
    \draw[->] (k) -- (attn);
    \draw[->] (v) -- (attn);

    \node[rectangle,draw,fill=blue!30,minimum width=2cm,minimum height=0.8cm] (out) at (0, 5.5) {Output};
    \draw[->] (attn) -- (out);
\end{tikzpicture}
\end{center}
\end{frame}

% ============================================================
% SELF ATTENTION EXAMPLE
% ============================================================
\begin{frame}{Self-Attention Example 📝}
\textbf{Sentence: "The animal didn't cross the street because it was too tired"}

\vspace{1em}

\textbf{Question: What does "it" refer to?}

\vspace{1em}

\begin{center}
\begin{tikzpicture}[scale=0.7]
    % Words
    \node at (0, 2) {The};
    \node at (1.2, 2) {animal};
    \node at (2.8, 2) {didn't};
    \node at (4, 2) {cross};
    \node at (5.2, 2) {the};
    \node at (6.2, 2) {street};
    \node at (7.5, 2) {because};
    \node[rectangle,draw,fill=orange!30] at (8.8, 2) {it};
    \node at (9.8, 2) {was};
    \node at (10.8, 2) {too};
    \node at (11.8, 2) {tired};

    % Attention from "it"
    \draw[->, thick, line width=3pt, red!50] (8.8, 1.7) to[bend right=40] (1.2, 1.7);
    \node[red] at (5, 0.5) {\textbf{Strong attention to "animal"}};

    \draw[->, thick, line width=1pt, blue!30] (8.8, 1.7) to[bend right=30] (6.2, 1.7);
    \node[blue] at (8, 0) {\tiny weak attention to "street"};
\end{tikzpicture}
\end{center}

\vspace{1em}

\textbf{Self-attention allows the model to:}
\begin{itemize}
    \item Resolve pronouns
    \item Understand long-range dependencies
    \item Capture syntactic and semantic relationships
    \item Do this in parallel for all positions!
\end{itemize}
\end{frame}

% ============================================================
% MULTI HEAD ATTENTION
% ============================================================
\begin{frame}{Multi-Head Attention 🎭}
\textbf{Why use multiple attention heads?}

\vspace{1em}

\begin{columns}[T]
\begin{column}{0.48\textwidth}
\textbf{Intuition:}
\begin{itemize}
    \item Different heads learn different relationships
    \item Head 1: Syntactic relationships
    \item Head 2: Semantic relationships
    \item Head 3: Positional patterns
    \item Etc.
\end{itemize}

\vspace{1em}

\textbf{Formula:}
\begin{align*}
\text{MultiHead}(Q,K,V) &= \text{Concat}(h_1, ..., h_H)W_O \\
h_i &= \text{Attention}(QW_i^Q, KW_i^K, VW_i^V)
\end{align*}
\end{column}

\begin{column}{0.48\textwidth}
\begin{center}
\begin{tikzpicture}[scale=0.65]
    \node[rectangle,draw,fill=blue!20] (input) at (2, 0) {Input};

    % Multiple heads
    \foreach \x/\name in {0/Head 1, 1.5/Head 2, 3/Head 3, 4.5/Head H} {
        \node[rectangle,draw,fill=green!20,minimum width=1cm,minimum height=0.6cm] (h\x) at (\x, 2) {\tiny\name};
        \draw[->] (input) -- (h\x);
    }

    \node at (5.5, 2) {\tiny ...};

    % Concatenate
    \node[rectangle,draw,fill=orange!20,minimum width=5cm,minimum height=0.8cm] (concat) at (2, 4) {Concatenate};
    \foreach \x in {0, 1.5, 3, 4.5} {
        \draw[->] (h\x) -- (concat);
    }

    % Linear
    \node[rectangle,draw,fill=purple!20,minimum width=3cm,minimum height=0.8cm] (linear) at (2, 5.5) {Linear ($W_O$)};
    \draw[->] (concat) -- (linear);

    % Output
    \node[rectangle,draw,fill=red!20] (output) at (2, 7) {Output};
    \draw[->] (linear) -- (output);
\end{tikzpicture}
\end{center}

\vspace{0.5em}
\textbf{Typical:} 8-16 heads in practice
\end{column}
\end{columns}

\vspace{0.5em}
\textit{BERT-base: 12 heads, BERT-large: 16 heads, GPT-3: 96 heads!}
\end{frame}

% ============================================================
% TYPES OF ATTENTION
% ============================================================
\begin{frame}{Three Types of Attention 🔍}
\begin{enumerate}
    \item \textbf{Self-Attention (Encoder)}
    \begin{itemize}
        \item Each position attends to all positions in same sequence
        \item Bidirectional: can see past and future
        \item Used in: BERT, encoder-only models
    \end{itemize}

    \vspace{0.5em}

    \item \textbf{Masked Self-Attention (Decoder)}
    \begin{itemize}
        \item Each position attends only to previous positions
        \item Prevents "looking into the future"
        \item Used in: GPT, decoder-only models
    \end{itemize}

    \vspace{0.5em}

    \item \textbf{Cross-Attention (Encoder-Decoder)}
    \begin{itemize}
        \item Decoder attends to encoder outputs
        \item Queries from decoder, Keys/Values from encoder
        \item Used in: T5, BART, machine translation
    \end{itemize}
\end{enumerate}

\vspace{1em}

\begin{center}
\begin{tabular}{lcc}
\toprule
\textbf{Type} & \textbf{Q from} & \textbf{K, V from} \\
\midrule
Self-Attention & Sequence & Same sequence \\
Masked Self-Attention & Sequence & Same sequence (past only) \\
Cross-Attention & Decoder & Encoder \\
\bottomrule
\end{tabular}
\end{center}
\end{frame}

% ============================================================
% POSITIONAL ENCODING
% ============================================================
\begin{frame}{Positional Encoding 📍}
\textbf{Problem: Self-attention is permutation-invariant!}

\vspace{0.5em}

Without position information:
\begin{itemize}
    \item "The cat sat" = "sat cat the" = "cat the sat"
    \item Order matters in language!
\end{itemize}

\vspace{1em}

\textbf{Solution: Add positional encodings to input embeddings}

\vspace{1em}

\begin{columns}[T]
\begin{column}{0.48\textwidth}
\textbf{Original Transformer (Sinusoidal):}
\begin{align*}
PE_{(pos, 2i)} &= \sin(pos / 10000^{2i/d}) \\
PE_{(pos, 2i+1)} &= \cos(pos / 10000^{2i/d})
\end{align*}

\textbf{Properties:}
\begin{itemize}
    \item Deterministic
    \item Unique for each position
    \item Smooth changes
    \item Generalizes to unseen lengths
\end{itemize}
\end{column}

\begin{column}{0.48\textwidth}
\textbf{Modern Approach: RoPE}

Rotary Position Embedding (Su et al. 2021)
\begin{itemize}
    \item Rotates Q and K in complex space
    \item Relative position encoding
    \item Better extrapolation
    \item Used in: LLaMA, GPT-NeoX
\end{itemize}

\vspace{1em}

\textbf{Learned Positional Embeddings:}
\begin{itemize}
    \item Treat positions as vocabulary
    \item Learn embeddings during training
    \item Used in: BERT, GPT
\end{itemize}
\end{column}
\end{columns}

\vspace{0.5em}
\small
\textit{References: Vaswani et al. (2017), Su et al. (2021) - "RoFormer: Enhanced Transformer with Rotary Position Embedding"}
\end{frame}

% ============================================================
% FEED FORWARD
% ============================================================
\begin{frame}[fragile]{Feed-Forward Networks 🔄}
\textbf{After attention, apply position-wise feed-forward network}

\vspace{1em}

\textbf{Architecture:}
\begin{align*}
\text{FFN}(x) = \max(0, xW_1 + b_1)W_2 + b_2
\end{align*}

\begin{itemize}
    \item Two linear transformations with ReLU activation
    \item Applied to each position independently
    \item Same network for all positions
    \item Typical: Expand 4x then project back
\end{itemize}

\vspace{1em}

\begin{center}
\begin{tikzpicture}[scale=0.8]
    \node[rectangle,draw,fill=blue!20,minimum width=2cm,minimum height=0.8cm,align=center] (input) at (0, 0) {Input\\(768 dims)};

    \node[rectangle,draw,fill=green!20,minimum width=2cm,minimum height=0.8cm,align=center] (expand) at (0, 2) {Linear + ReLU\\(3072 dims)};
    \draw[->] (input) -- (expand);

    \node[rectangle,draw,fill=orange!20,minimum width=2cm,minimum height=0.8cm,align=center] (project) at (0, 4) {Linear\\(768 dims)};
    \draw[->] (expand) -- (project);

    \node[rectangle,draw,fill=red!20,minimum width=2cm,minimum height=0.8cm,align=center] (output) at (0, 6) {Output\\(768 dims)};
    \draw[->] (project) -- (output);

    \node[right=0.5cm of expand] {4x expansion};
\end{tikzpicture}
\end{center}

\textbf{Purpose:} Add non-linearity and transform representations
\end{frame}

% ============================================================
% LAYER NORM AND RESIDUALS
% ============================================================
\begin{frame}{Layer Normalization \& Residual Connections 🔗}
\textbf{Critical for training deep transformers!}

\vspace{1em}

\begin{columns}[T]
\begin{column}{0.48\textwidth}
\textbf{Residual Connections:}
\begin{align*}
\text{Output} = \text{Layer}(x) + x
\end{align*}

\begin{itemize}
    \item Allows gradients to flow directly
    \item Prevents vanishing gradients
    \item Enables very deep networks
    \item Identity mapping as fallback
\end{itemize}

\vspace{0.5em}

\textbf{Analogy:}
Like highway roads for information flow!
\end{column}

\begin{column}{0.48\textwidth}
\textbf{Layer Normalization:}
\begin{align*}
\text{LayerNorm}(x) = \gamma \frac{x - \mu}{\sigma} + \beta
\end{align*}

\begin{itemize}
    \item Normalize across features
    \item Stabilizes training
    \item Faster convergence
    \item Independent of batch size
\end{itemize}

\vspace{0.5em}

\begin{center}
\begin{tikzpicture}[scale=0.6]
    \node[rectangle,draw] (x) at (0, 0) {$x$};
    \node[rectangle,draw] (layer) at (0, 1.5) {Layer};
    \node[rectangle,draw] (ln) at (0, 3) {LayerNorm};
    \node[rectangle,draw] (add) at (0, 4.5) {Add};

    \draw[->] (x) -- (layer);
    \draw[->] (layer) -- (ln);
    \draw[->] (ln) -- (add);
    \draw[->, thick, red] (x) to[bend left=40] (add);

    \node[red,right=0.2cm of add] {Residual};
\end{tikzpicture}
\end{center}
\end{column}
\end{columns}

\vspace{1em}

\textbf{Standard Pattern:}
\begin{align*}
x &\leftarrow \text{LayerNorm}(x + \text{MultiHeadAttention}(x)) \\
x &\leftarrow \text{LayerNorm}(x + \text{FFN}(x))
\end{align*}
\end{frame}

% ============================================================
% FLASH ATTENTION
% ============================================================
\begin{frame}{FlashAttention: Making Transformers Faster ⚡}
\textbf{Problem: Standard attention is slow and memory-hungry!}

\vspace{0.5em}

\begin{alertblock}{Standard Attention Complexity}
\begin{itemize}
    \item Time: $O(n^2)$ where $n$ is sequence length
    \item Memory: $O(n^2)$ to store attention matrix
    \item Bottleneck: Reading/writing to GPU memory (HBM)
\end{itemize}
\end{alertblock}

\vspace{1em}

\begin{columns}[T]
\begin{column}{0.48\textwidth}
\textbf{FlashAttention Innovation:}
\begin{itemize}
    \item Tile-based computation
    \item Uses fast SRAM instead of slow HBM
    \item Fused operations
    \item Recomputation in backward pass
    \item \alert{2-4x faster!}
    \item Enables longer sequences
\end{itemize}
\end{column}

\begin{column}{0.48\textwidth}
\textbf{Impact:}
\begin{itemize}
    \item GPT-3: 512 → 2048 context
    \item BERT: Faster training
    \item Long-document understanding
    \item Used in: LLaMA, GPT-4
\end{itemize}

\vspace{1em}

\begin{center}
\begin{tabular}{lcc}
\toprule
\textbf{Method} & \textbf{Speed} & \textbf{Memory} \\
\midrule
Standard & 1x & 1x \\
FlashAttn & \textbf{3x} & \textbf{0.5x} \\
\bottomrule
\end{tabular}
\end{center}
\end{column}
\end{columns}

\vspace{0.5em}
\small
\textit{Reference: Dao et al. (2022) - "FlashAttention: Fast and Memory-Efficient Exact Attention"}
\end{frame}

% ============================================================
% ENCODER VS DECODER VS BOTH
% ============================================================
\begin{frame}{Three Transformer Architectures 🏗️}
\begin{center}
\begin{tabular}{p{3cm}p{4cm}p{4cm}}
\toprule
\textbf{Architecture} & \textbf{How it works} & \textbf{Examples \& Uses} \\
\midrule
\textbf{Encoder-only} &
Bidirectional self-attention. Can see entire sequence. &
BERT, RoBERTa, ALBERT\\
\textit{Best for:} Classification, NER, QA \\
\midrule
\textbf{Decoder-only} &
Masked self-attention. Can only see past. Autoregressive. &
GPT, GPT-2, GPT-3, LLaMA\\
\textit{Best for:} Generation, completion \\
\midrule
\textbf{Encoder-Decoder} &
Encoder: bidirectional\\
Decoder: masked + cross-attention &
T5, BART, mT5\\
\textit{Best for:} Translation, summarization \\
\bottomrule
\end{tabular}
\end{center}

\vspace{1em}

\begin{block}{Key Difference: Attention Masking}
\begin{itemize}
    \item \textbf{Encoder:} Full self-attention (bidirectional)
    \item \textbf{Decoder:} Causal/masked attention (unidirectional)
\end{itemize}
\end{block}
\end{frame}

% ============================================================
% VISUAL ARCHITECTURES
% ============================================================
\begin{frame}{Visual Comparison: Encoder vs Decoder vs Both}
\begin{center}
\begin{tikzpicture}[scale=0.5]
    % Encoder only (BERT)
    \node at (0, 6) {\textbf{Encoder-Only (BERT)}};
    \foreach \y in {0, 1.5, 3} {
        \node[rectangle,draw,fill=blue!30,minimum width=3cm,minimum height=1cm] at (0, \y) {Self-Attention};
    }
    \draw[<->, thick] (-1.8, 1.5) -- (-1.8, 4);
    \node[left] at (-2, 2.7) {\tiny bidirectional};

    % Decoder only (GPT)
    \node at (6, 6) {\textbf{Decoder-Only (GPT)}};
    \foreach \y in {0, 1.5, 3} {
        \node[rectangle,draw,fill=red!30,minimum width=3cm,minimum height=1cm] at (6, \y) {Masked Attn};
    }
    \draw[->, thick] (4.2, 0.5) -- (4.2, 3.5);
    \node[left] at (3.8, 2) {\tiny causal};

    % Encoder-Decoder (T5)
    \node at (13, 6) {\textbf{Encoder-Decoder (T5)}};
    \node[rectangle,draw,fill=blue!30,minimum width=2cm,minimum height=2.5cm] at (11, 1.5) {Encoder};
    \node[rectangle,draw,fill=red!30,minimum width=2cm,minimum height=2.5cm] at (15, 1.5) {Decoder};
    \draw[->, thick] (12.2, 1.5) -- (13.8, 1.5) node[midway,above] {\tiny cross-attn};
\end{tikzpicture}
\end{center}

\vspace{1em}

\textbf{Choosing the Right Architecture:}
\begin{itemize}
    \item Need to understand full context? → \textbf{Encoder} (BERT)
    \item Need to generate text? → \textbf{Decoder} (GPT)
    \item Need both (translate, summarize)? → \textbf{Encoder-Decoder} (T5)
\end{itemize}
\end{frame}

% ============================================================
% TRANSFORMER CODE
% ============================================================
\begin{frame}[fragile]{Using Transformers in Practice 💻}
\textbf{HuggingFace makes it easy!}

\begin{lstlisting}[language=Python]
from transformers import AutoModel, AutoTokenizer

# Load pre-trained model
model_name = "bert-base-uncased"
tokenizer = AutoTokenizer.from_pretrained(model_name)
model = AutoModel.from_pretrained(model_name)

# Tokenize input
text = "The animal didn't cross the street because it was too tired"
inputs = tokenizer(text, return_tensors="pt")

# Get contextualized embeddings
outputs = model(**inputs)
hidden_states = outputs.last_hidden_state  # Shape: [batch, seq_len, 768]

# Extract embedding for specific token
token_embeddings = hidden_states[0]  # [seq_len, 768]
print(f"Shape: {token_embeddings.shape}")
# Each token now has a context-aware representation!
\end{lstlisting}

\vspace{0.5em}
\small
\textit{Reference: HuggingFace Course - Chapter 1.4 (How do Transformers work?)}
\end{frame}

% ============================================================
% BERT INTRO
% ============================================================
\begin{frame}{BERT: Bidirectional Encoder Representations 🎯}
\begin{center}
\LARGE
\textbf{BERT = Encoder-only Transformer}
\end{center}

\vspace{1em}

\textbf{Key Innovation: Masked Language Modeling (MLM)}

\vspace{1em}

\begin{columns}[T]
\begin{column}{0.48\textwidth}
\textbf{Traditional Language Models:}
\begin{itemize}
    \item Left-to-right (GPT)
    \item Or right-to-left
    \item \alert{Unidirectional context}
    \item Can't see full picture
\end{itemize}

\vspace{1em}

\textbf{Example:}
\begin{itemize}
    \item "The cat sat on the \_\_\_"
    \item Only sees left context
\end{itemize}
\end{column}

\pause

\begin{column}{0.48\textwidth}
\textbf{BERT (MLM):}
\begin{itemize}
    \item Mask random tokens
    \item Predict from \alert{bidirectional context}
    \item Sees both left \& right
    \item Deeper understanding
\end{itemize}

\vspace{1em}

\textbf{Example:}
\begin{itemize}
    \item "The cat [MASK] on the mat"
    \item Sees: "The cat" AND "on the mat"
    \item Predicts: "sat"
\end{itemize}
\end{column}
\end{columns}

\vspace{1em}
\small
\textit{Reference: Devlin et al. (2019) - "BERT: Pre-training of Deep Bidirectional Transformers"}
\end{frame}

% ============================================================
% MLM TRAINING
% ============================================================
\begin{frame}{Masked Language Modeling (MLM) 🎭}
\textbf{BERT's Pre-training Objective}

\vspace{1em}

\textbf{Training Procedure:}
\begin{enumerate}
    \item Take a sentence
    \item Randomly mask 15\% of tokens
    \item Of the masked tokens:
    \begin{itemize}
        \item 80\%: Replace with [MASK]
        \item 10\%: Replace with random word
        \item 10\%: Keep unchanged
    \end{itemize}
    \item Predict the original tokens
\end{enumerate}

\vspace{1em}

\begin{exampleblock}{Example}
\textbf{Original:} "My dog is hairy"

\textbf{Masked:} "My dog is [MASK]"

\textbf{Labels:} [-, -, -, hairy]

\textbf{Prediction:} Model predicts "hairy" using bidirectional context
\end{exampleblock}

\vspace{1em}

\textbf{Why the 80/10/10 split?}
\begin{itemize}
    \item Prevents overfitting to [MASK] token
    \item Forces model to maintain representations for all tokens
\end{itemize}
\end{frame}

% ============================================================
% BERT ARCHITECTURE
% ============================================================
\begin{frame}{BERT Architecture Variants 📊}
\begin{center}
\begin{tabular}{lcccc}
\toprule
\textbf{Model} & \textbf{Layers} & \textbf{Hidden Size} & \textbf{Heads} & \textbf{Params} \\
\midrule
BERT-Base & 12 & 768 & 12 & 110M \\
BERT-Large & 24 & 1024 & 16 & 340M \\
RoBERTa-Base & 12 & 768 & 12 & 125M \\
RoBERTa-Large & 24 & 1024 & 16 & 355M \\
ALBERT-Base & 12 & 768 & 12 & 12M \\
DistilBERT & 6 & 768 & 12 & 66M \\
\bottomrule
\end{tabular}
\end{center}

\vspace{1em}

\textbf{Architecture Details (BERT-Base):}
\begin{itemize}
    \item 12 transformer encoder layers
    \item 768-dimensional hidden states
    \item 12 attention heads per layer
    \item 3072-dimensional feed-forward intermediate size
    \item Maximum sequence length: 512 tokens
    \item Vocabulary size: 30,000 WordPiece tokens
\end{itemize}
\end{frame}

% ============================================================
% BERT VARIANTS
% ============================================================
\begin{frame}{BERT Variants \& Improvements 🔬}
\begin{enumerate}
    \item \textbf{RoBERTa (Liu et al. 2019)}
    \begin{itemize}
        \item Remove Next Sentence Prediction (NSP)
        \item Train longer with bigger batches
        \item Dynamic masking
        \item \alert{Better performance than BERT}
    \end{itemize}

    \vspace{0.5em}

    \item \textbf{ALBERT (Lan et al. 2019)}
    \begin{itemize}
        \item Parameter sharing across layers
        \item Factorized embedding parameterization
        \item \alert{12M params vs BERT's 110M}
        \item Similar performance, much smaller
    \end{itemize}

    \vspace{0.5em}

    \item \textbf{DistilBERT (Sanh et al. 2019)}
    \begin{itemize}
        \item Knowledge distillation from BERT
        \item 6 layers instead of 12
        \item \alert{40\% smaller, 60\% faster}
        \item Retains 97\% of performance
    \end{itemize}
\end{enumerate}

\vspace{0.5em}
\small
\textit{References: Liu et al. (2019) - RoBERTa, Sanh et al. (2019) - DistilBERT}
\end{frame}

% ============================================================
% BERT FINE-TUNING
% ============================================================
\begin{frame}[fragile]{Fine-tuning BERT 🎓}
\textbf{Two-stage process: Pre-train then Fine-tune}

\vspace{1em}

\begin{lstlisting}[language=Python]
from transformers import BertForSequenceClassification, Trainer, TrainingArguments

# Load pre-trained BERT with classification head
model = BertForSequenceClassification.from_pretrained(
    'bert-base-uncased',
    num_labels=2  # Binary classification
)

# Define training arguments
training_args = TrainingArguments(
    output_dir='./results',
    num_train_epochs=3,
    per_device_train_batch_size=16,
    learning_rate=2e-5,
    warmup_steps=500,
)

# Train
trainer = Trainer(
    model=model,
    args=training_args,
    train_dataset=train_dataset,
    eval_dataset=eval_dataset,
)

trainer.train()
\end{lstlisting}

\vspace{0.5em}
\small
\textit{Reference: HuggingFace Course - Chapters 1.5 (Using Transformers) and 7.3 (Fine-tuning)}
\end{frame}

% ============================================================
% BERT APPLICATIONS
% ============================================================
\begin{frame}{BERT Applications 🚀}
\textbf{BERT excels at understanding tasks:}

\vspace{1em}

\begin{columns}[T]
\begin{column}{0.48\textwidth}
\textbf{Classification Tasks:}
\begin{itemize}
    \item Sentiment Analysis
    \item Topic Classification
    \item Spam Detection
    \item Intent Recognition
\end{itemize}

\vspace{0.5em}

\textbf{Token-Level Tasks:}
\begin{itemize}
    \item Named Entity Recognition (NER)
    \item Part-of-Speech Tagging
    \item Word Sense Disambiguation
\end{itemize}
\end{column}

\begin{column}{0.48\textwidth}
\textbf{Span-Level Tasks:}
\begin{itemize}
    \item Question Answering
    \item Extractive Summarization
    \item Information Extraction
\end{itemize}

\vspace{0.5em}

\textbf{Sentence-Pair Tasks:}
\begin{itemize}
    \item Semantic Similarity
    \item Natural Language Inference
    \item Paraphrase Detection
\end{itemize}
\end{column}
\end{columns}

\vspace{1em}

\begin{block}{Industry Impact}
BERT powers:
\begin{itemize}
    \item Google Search (understanding queries)
    \item Customer service chatbots
    \item Content moderation
    \item Document understanding
\end{itemize}
\end{block}
\end{frame}

% ============================================================
% CONTEXTUAL EMBEDDINGS EXAMPLE
% ============================================================
\begin{frame}[fragile]{Contextual Embeddings in Action 🔬}
\textbf{Remember "bank"? Let's see BERT handle it!}

\begin{lstlisting}[language=Python]
from transformers import BertTokenizer, BertModel
import torch

tokenizer = BertTokenizer.from_pretrained('bert-base-uncased')
model = BertModel.from_pretrained('bert-base-uncased')

# Two different contexts for "bank"
sent1 = "I deposited money at the bank"
sent2 = "We sat by the river bank"

# Get embeddings
def get_embedding(sentence, target_word):
    inputs = tokenizer(sentence, return_tensors='pt')
    outputs = model(**inputs)
    # Find position of target word
    tokens = tokenizer.tokenize(sentence)
    idx = tokens.index(target_word) + 1  # +1 for [CLS]
    return outputs.last_hidden_state[0, idx, :]

emb1 = get_embedding(sent1, "bank")  # Financial bank
emb2 = get_embedding(sent2, "bank")  # River bank

# Compare similarity
similarity = torch.cosine_similarity(emb1, emb2, dim=0)
print(f"Similarity: {similarity:.3f}")  # Low! (~0.3-0.5)
# Different contexts → Different embeddings!
\end{lstlisting}
\end{frame}

% ============================================================
% COGNITIVE NEUROSCIENCE
% ============================================================
\begin{frame}{Cognitive Neuroscience Perspective 🧠}
\textbf{How do brains and models process language?}

\vspace{1em}

\begin{columns}[T]
\begin{column}{0.48\textwidth}
\textbf{Predictive Processing in the Brain:}
\begin{itemize}
    \item Brain constantly predicts upcoming input
    \item N400: Neural response to unexpected words
    \item P600: Syntactic anomaly detection
    \item Context shapes predictions
    \item Prediction errors drive learning
\end{itemize}

\vspace{0.5em}

\textbf{Key Brain Regions:}
\begin{itemize}
    \item \textbf{Left IFG}: Syntax processing
    \item \textbf{Left STG/MTG}: Semantic processing
    \item \textbf{ATL}: Conceptual knowledge
\end{itemize}
\end{column}

\begin{column}{0.48\textwidth}
\textbf{Predictive Processing in Models:}
\begin{itemize}
    \item BERT: Predict masked words
    \item GPT: Predict next word
    \item Both use context to predict
    \item Surprise = high loss
    \item Gradient descent = learning
\end{itemize}

\vspace{0.5em}

\textbf{Similarities:}
\begin{itemize}
    \item Both hierarchical
    \item Both context-sensitive
    \item Both predictive
    \item Both learn from errors
\end{itemize}
\end{column}
\end{columns}

\vspace{0.5em}
\small
\textit{References: Kuperberg \& Jaeger (2016), Willems et al. (2016), Hagoort \& Indefrey (2014)}
\end{frame}

% ============================================================
% PREDICTION IN BRAINS VS MODELS
% ============================================================
\begin{frame}{Prediction in Brains vs. Language Models 🧠🤖}
\textbf{Parallels between neural and artificial systems}

\vspace{1em}

\begin{center}
\begin{tabular}{p{5cm}p{5cm}}
\toprule
\textbf{Human Brain} & \textbf{Transformer Models} \\
\midrule
N400 component\\
(surprise at unexpected word) &
High cross-entropy loss\\
(low probability prediction) \\
\midrule
Hierarchical processing\\
(sounds → words → sentences) &
Layer-by-layer representations\\
(tokens → phrases → meaning) \\
\midrule
Context integration\\
(prior discourse, world knowledge) &
Self-attention mechanism\\
(attend to relevant context) \\
\midrule
Prediction error signals &
Gradient-based learning \\
\midrule
Distributed representations\\
(population coding) &
Distributed embeddings\\
(vector representations) \\
\bottomrule
\end{tabular}
\end{center}

\vspace{1em}

\textbf{Question:} Are these superficial analogies or deep connections?

\small
\textit{Reference: Kuperberg \& Jaeger (2016) - "What do we mean by prediction in language comprehension?"}
\end{frame}

% ============================================================
% NEURAL ENCODING
% ============================================================
\begin{frame}{Neural Encoding with Language Models 🔬}
\textbf{Can we predict brain activity from language models?}

\vspace{1em}

\textbf{Approach:}
\begin{enumerate}
    \item Show participants text while recording brain activity (fMRI/EEG)
    \item Extract representations from language model (e.g., BERT layer 6)
    \item Train regression model: Model embeddings → Brain activity
    \item Test: Can we predict brain responses to new text?
\end{enumerate}

\vspace{1em}

\textbf{Findings:}
\begin{itemize}
    \item \textbf{Better models → Better brain prediction}
    \item BERT outperforms Word2Vec at predicting brain activity
    \item Different layers correlate with different brain regions
    \item Middle layers best for semantic areas
    \item Suggests shared representations?
\end{itemize}

\vspace{1em}

\begin{alertblock}{But...}
Correlation ≠ Causation! Models might capture statistical regularities that brains use, without using the same mechanisms.
\end{alertblock}

\small
\textit{Reference: Willems et al. (2016) - "Prediction during natural language comprehension"}
\end{frame}

% ============================================================
% WHAT DOES MODEL UNDERSTAND
% ============================================================
\begin{frame}{Discussion: What Does the Model "Understand"? 🤔}
\begin{center}
\LARGE
\textbf{Does BERT understand language?}
\end{center}

\vspace{1em}

\begin{columns}[T]
\begin{column}{0.48\textwidth}
\textbf{Evidence FOR understanding:}
\begin{itemize}
    \item Captures syntax and semantics
    \item Resolves ambiguity
    \item Handles long-range dependencies
    \item Generalizes to new examples
    \item Predicts brain activity
    \item Solves complex tasks
\end{itemize}

\vspace{0.5em}
\textit{"If it acts like it understands, maybe it does?"}
\end{column}

\begin{column}{0.48\textwidth}
\textbf{Evidence AGAINST understanding:}
\begin{itemize}
    \item No grounding in physical world
    \item No sensory experience
    \item No social context
    \item Brittle to adversarial examples
    \item No common sense reasoning
    \item Only pattern matching?
\end{itemize}

\vspace{0.5em}
\textit{"Understanding requires more than statistical patterns"}
\end{column}
\end{columns}

\vspace{1em}

\begin{block}{Key Questions}
\begin{itemize}
    \item What is the difference between understanding and correlation?
    \item Can meaning exist without grounding?
    \item Is human understanding fundamentally different?
\end{itemize}
\end{block}

\textbf{Class Discussion: What do YOU think?}
\end{frame}

% ============================================================
% LIMITATIONS
% ============================================================
\begin{frame}{Limitations of Current Models ⚠️}
\textbf{Despite impressive performance, transformers have limitations:}

\vspace{1em}

\begin{enumerate}
    \item \textbf{Quadratic Complexity}
    \begin{itemize}
        \item Self-attention scales as $O(n^2)$
        \item Limited context windows (512-4096 tokens)
        \item Cannot process very long documents efficiently
    \end{itemize}

    \item \textbf{No True Understanding}
    \begin{itemize}
        \item Pattern matching vs. comprehension
        \item Lack of common sense
        \item No world model
    \end{itemize}

    \item \textbf{Data Efficiency}
    \begin{itemize}
        \item Requires massive training data
        \item Humans learn language with much less data
        \item Not biologically plausible
    \end{itemize}

    \item \textbf{Biases and Fairness}
    \begin{itemize}
        \item Inherits biases from training data
        \item Can amplify stereotypes
        \item Ethical concerns
    \end{itemize}
\end{enumerate}
\end{frame}

% ============================================================
% ANIMATED TRANSFORMER
% ============================================================
\begin{frame}{Interactive Learning Resources 🎓}
\textbf{Great tutorials for deeper understanding:}

\vspace{1em}

\begin{enumerate}
    \item \textbf{The Illustrated Transformer}
    \begin{itemize}
        \item Blog post by Jay Alammar
        \item Visual step-by-step walkthrough
        \item \url{https://jalammar.github.io/illustrated-transformer/}
    \end{itemize}

    \vspace{0.5em}

    \item \textbf{Animated Transformer Tutorial}
    \begin{itemize}
        \item Interactive visualizations
        \item Watch attention in action
        \item \url{https://transformer-circuits.pub/}
    \end{itemize}

    \vspace{0.5em}

    \item \textbf{HuggingFace Course}
    \begin{itemize}
        \item Chapter 1.4: How do Transformers work?
        \item Hands-on coding tutorials
        \item \url{https://huggingface.co/course/chapter1/4}
    \end{itemize}

    \vspace{0.5em}

    \item \textbf{Attention Mechanism Visualization}
    \begin{itemize}
        \item BertViz: Visualize BERT attention
        \item See which words attend to which
        \item \url{https://github.com/jessevig/bertviz}
    \end{itemize}
\end{enumerate}
\end{frame}

% ============================================================
% CODE EXAMPLE COMPARISON
% ============================================================
\begin{frame}[fragile]{Comparing Architectures: Code Examples 💻}
\textbf{Different architectures for different tasks}

\begin{lstlisting}[language=Python]
from transformers import AutoModelForSequenceClassification, \
                         AutoModelForCausalLM, \
                         AutoModelForSeq2SeqLM

# 1. ENCODER-ONLY (BERT): Classification
encoder_model = AutoModelForSequenceClassification.from_pretrained(
    "bert-base-uncased", num_labels=2
)
# Use for: Sentiment analysis, NER, classification

# 2. DECODER-ONLY (GPT): Text generation
decoder_model = AutoModelForCausalLM.from_pretrained("gpt2")
# Use for: Text completion, creative writing, few-shot learning

# 3. ENCODER-DECODER (T5): Seq2Seq tasks
seq2seq_model = AutoModelForSeq2SeqLM.from_pretrained("t5-base")
# Use for: Translation, summarization, question answering

# All use transformer architecture, different attention patterns!
\end{lstlisting}

\vspace{0.5em}
\textbf{Key Takeaway:} Choose architecture based on your task!
\end{frame}

% ============================================================
% ASSIGNMENT 4
% ============================================================
\begin{frame}{Assignment 4: Context-Aware Models 📝}
\textbf{Hands-on experience with transformers!}

\vspace{1em}

\textbf{Tasks:}
\begin{enumerate}
    \item \textbf{Implement Attention Mechanism}
    \begin{itemize}
        \item Build scaled dot-product attention from scratch
        \item Visualize attention weights
    \end{itemize}

    \item \textbf{Fine-tune BERT}
    \begin{itemize}
        \item Load pre-trained BERT
        \item Fine-tune on sentiment analysis
        \item Compare to baseline models
    \end{itemize}

    \item \textbf{Analyze Contextual Embeddings}
    \begin{itemize}
        \item Extract embeddings for polysemous words
        \item Visualize how context changes representations
        \item Compare BERT vs Word2Vec
    \end{itemize}

    \item \textbf{Explore Different Architectures}
    \begin{itemize}
        \item Compare BERT (encoder) vs GPT (decoder)
        \item Test on different tasks
        \item Analyze strengths/weaknesses
    \end{itemize}

    \item \textbf{Research Component}
    \begin{itemize}
        \item Read one paper from references
        \item Write brief summary \& critical analysis
    \end{itemize}
\end{enumerate}

\vspace{0.5em}
\textbf{Due:} Check course website for deadline
\end{frame}

% ============================================================
% PRACTICAL TIPS
% ============================================================
\begin{frame}{Practical Tips for Working with Transformers 💡}
\begin{enumerate}
    \item \textbf{Start with Pre-trained Models}
    \begin{itemize}
        \item Don't train from scratch (too expensive!)
        \item Use HuggingFace Model Hub
        \item Choose appropriate model size
    \end{itemize}

    \item \textbf{Fine-tuning Best Practices}
    \begin{itemize}
        \item Use small learning rate (1e-5 to 5e-5)
        \item Add warmup steps
        \item Monitor for overfitting
        \item Freeze early layers if data is limited
    \end{itemize}

    \item \textbf{Computational Efficiency}
    \begin{itemize}
        \item Use mixed precision training (FP16)
        \item Gradient accumulation for larger batch sizes
        \item Consider DistilBERT for faster inference
        \item Use FlashAttention when available
    \end{itemize}

    \item \textbf{Evaluation}
    \begin{itemize}
        \item Use task-specific metrics
        \item Test on out-of-distribution data
        \item Check for biases
        \item Visualize attention for interpretability
    \end{itemize}
\end{enumerate}
\end{frame}

% ============================================================
% RESOURCES
% ============================================================
\begin{frame}{Resources \& Further Reading 📚}
\textbf{Key Papers:}
\begin{itemize}
    \item Sutskever et al. (2014) - Sequence to Sequence Learning
    \item Bahdanau et al. (2015) - Neural Machine Translation by Jointly Learning to Align
    \item Vaswani et al. (2017) - Attention Is All You Need
    \item Devlin et al. (2019) - BERT: Pre-training of Deep Bidirectional Transformers
    \item Liu et al. (2019) - RoBERTa: A Robustly Optimized BERT Pretraining Approach
    \item Sanh et al. (2019) - DistilBERT: A distilled version of BERT
    \item Su et al. (2021) - RoFormer: Enhanced Transformer with Rotary Position Embedding
    \item Dao et al. (2022) - FlashAttention: Fast and Memory-Efficient Exact Attention
\end{itemize}

\vspace{0.5em}

\textbf{Cognitive Neuroscience:}
\begin{itemize}
    \item Hagoort \& Indefrey (2014) - The neurobiology of language beyond single words
    \item Kuperberg \& Jaeger (2016) - What do we mean by prediction in language comprehension?
    \item Willems et al. (2016) - Prediction during natural language comprehension
\end{itemize}

\vspace{0.5em}

\textbf{Tutorials:}
\begin{itemize}
    \item HuggingFace Course: Chapters 1.4, 1.5, 7.3
    \item The Illustrated Transformer: \url{https://jalammar.github.io/illustrated-transformer/}
    \item Animated Transformer: \url{https://transformer-circuits.pub/}
\end{itemize}
\end{frame}

% ============================================================
% SUMMARY
% ============================================================
\begin{frame}{Summary 🎯}
\textbf{What we learned:}

\begin{enumerate}
    \item \textbf{Evolution of Context}
    \begin{itemize}
        \item Seq2Seq → Attention → Transformers
        \item From bottleneck to full parallelization
    \end{itemize}

    \item \textbf{Attention Mechanisms}
    \begin{itemize}
        \item Self-attention, masked attention, cross-attention
        \item Query, Key, Value framework
        \item Multi-head attention for diverse relationships
    \end{itemize}

    \item \textbf{Transformer Architecture}
    \begin{itemize}
        \item Pure attention-based architecture
        \item Encoder-only (BERT), Decoder-only (GPT), Both (T5)
        \item Positional encoding, layer norm, residual connections
    \end{itemize}

    \item \textbf{BERT \& Masked Language Modeling}
    \begin{itemize}
        \item Bidirectional context from MLM objective
        \item Pre-train then fine-tune paradigm
        \item Variants: RoBERTa, ALBERT, DistilBERT
    \end{itemize}

    \item \textbf{Brain-Model Parallels}
    \begin{itemize}
        \item Both use prediction
        \item Both hierarchical and context-sensitive
        \item Deep questions about understanding
    \end{itemize}
\end{enumerate}
\end{frame}

% ============================================================
% LOOKING FORWARD
% ============================================================
\begin{frame}{Looking Forward 🔮}
\textbf{Where do we go from here?}

\vspace{1em}

\textbf{Next Topics:}
\begin{itemize}
    \item Scaling laws: Bigger models, better performance?
    \item Efficient transformers: Linear attention, sparse attention
    \item Multimodal transformers: Vision + Language
    \item In-context learning: GPT-3 and beyond
    \item Alignment: Making models helpful and safe
\end{itemize}

\vspace{1em}

\textbf{Open Questions:}
\begin{itemize}
    \item How can we make transformers more efficient?
    \item Can we combine symbolic reasoning with neural networks?
    \item How do we ensure fairness and reduce bias?
    \item What is the role of grounding in understanding?
    \item Are we building towards AGI, or is something missing?
\end{itemize}

\vspace{1em}

\begin{center}
\Large
\textbf{The transformer revolution continues! 🚀}
\end{center}
\end{frame}

% ============================================================
% QUESTIONS
% ============================================================
\begin{frame}{Questions? 🙋}
\begin{center}
\LARGE
\textbf{Discussion Time}
\end{center}

\vspace{2em}

\textbf{Topics to discuss:}
\begin{itemize}
    \item Attention mechanisms
    \item Transformer architectures
    \item BERT vs GPT: When to use which?
    \item Do models really "understand"?
    \item Assignment 4 questions
\end{itemize}

\vspace{2em}

\begin{center}
\Large
Thank you! 🙏
\end{center}
\end{frame}

\end{document}
